\section{Introduction}
\label{sec:intro}
% src: no result is previewed here (STYLE rule 72). Citations: refs.bib and refs_extra.bib entries only.
Retrieval-augmented generation (RAG) lets a language model answer from documents that it was not trained on, by adding retrieved passages to its input \citep{lewis2020rag,guu2020realm,ram2023incontext}. The idea grew from open-domain question answering, in which a retriever finds passages and a reader extracts or writes the answer \citep{chen2017drqa,lee2019orqa,karpukhin2020dpr}. It now covers the fusion of many passages in the decoder, retrieval during pre-training and generation, and retrieval for models that are reached only through an interface \citep{izacard2021fid,borgeaud2022retro,izacard2023atlas,shi2024replug}. Surveys describe a quickly growing set of modules and variants \citep{gao2023ragsurvey,fan2024ragsurvey}. In practice, RAG is a common way to give a language model private, recent or domain-specific knowledge without retraining it. However, a RAG system answers well only if the right evidence reaches the reader, which is the language model that writes the answer. Errors in retrieval, in selection and in the length limit of the prompt can each keep evidence away from the reader, and the answer alone does not show which of them occurred \citep{ru2024ragchecker,joren2025sufficient}.

A RAG system has to make several decisions before the reader sees any text. The first is how to search. The options are a lexical function such as BM25 \citep{robertson2009bm25}, a dense encoder \citep{karpukhin2020dpr,xiong2021ance,izacard2022contriever}, a late-interaction or learned-sparse model \citep{khattab2020colbert,santhanam2022colbertv2,formal2021spladev2,lassance2024spladev3}, and a model that produces several signals at once \citep{chen2024bgem3,wang2024me5}. They can be used alone or together. Hybrid lexical and dense retrieval can raise recall \citep{kuzi2020hybrid}. The ranked lists are merged by rank \citep{cormack2009rrf} or by score \citep{bruch2023fusion,fox1993combination}. The second decision is what to do with the candidates. A cross-encoder \citep{nogueira2019monobert,nogueira2020monot5,lin2021pretrained} or a language model \citep{sun2023rankgpt,pradeep2023rankvicuna,pradeep2023rankzephyr,zhuang2024setwise} can reorder them, selection methods can choose a set of passages instead of a ranking \citep{lee2025setr,saxena2025meteora}, and filtering or compression can shorten what is kept \citep{wang2023filco,xu2024recomp,jiang2024longllmlingua}. The third decision concerns the question and the index, through query rewriting \citep{gao2023hyde,wang2023query2doc,ma2023queryrewriting,rackauckas2024ragfusion,zheng2024stepback}, the size and form of chunks \citep{chen2024densex,gunther2024latechunking}, and hierarchical or graph indexes \citep{sarthi2024raptor,edge2024graphrag}. The last decisions are how much text the reader receives, which is called the evidence budget, and whether to search once, search iteratively or read long documents instead. The text that the reader receives is called the evidence pack.

Persian and English differ in ways that matter for these decisions. Persian uses the Arabic script with zero-width non-joiners and several digit systems, so text is normalised before lexical matching \citep{mohtaj2018parsivar}. Tokenizers also split the same content into more tokens in some languages than in others \citep{petrov2023tokenizers,ahia2023cost}. Resources for Persian cover language understanding, question answering, retrieval and text embeddings \citep{farahani2021parsbert,khashabi2021parsinlu,darvishi2023pquad,aleahmad2009hamshahri,zinvandi2025famteb}, and they now include language models and RAG systems for Persian \citep{rostami2024persianmind,hosseini2024persianrag,hosseinbeigi2025persianrag}. Multilingual retrieval and RAG benchmarks list Persian among their languages \citep{zhang2023miracl,lawrie2023neuclir,thakur2025miragebench}, and cross-lingual retrieval and question answering resources support queries and documents in different languages \citep{zhang2021mrtydi,asai2021xorqa,asai2021cora,longpre2021mkqa,feng2022labse}. Studies of multilingual RAG report that quality varies with the language of the question, of the documents and of the reader \citep{chirkova2024mrag,wu2024multilingual,sharma2025fauxpolyglot,li2025culturalrag,liu2025xrag,amiraz2025crosslingual}. These findings suggest that conclusions from English benchmarks cannot be assumed to hold for Persian.

Existing comparisons address these decisions only in part. Studies that compare RAG setups to guide practice vary component choices or compare pipelines on English benchmarks \citep{wang2024bestpractices,rau2024bergen,li2024selfroute,ferrazzi2026agentic}. The closest study of tight budgets finds that a measure of whether the answer survives packing into the budget adds predictive information beyond recall \citep{bala2026recall}. A stress test of multi-hop RAG argues that budgets must be fixed across variants and that cluster-aware tests remove many apparent gains \citep{chacon2026fixedbudget}. Work on the evidence budget shows that readers use the middle of a long context poorly and that more passages can first help and then hurt \citep{liu2024lost,levy2025moredocs,jin2025longcontext}. Retrieval and long-context reading are complementary, since each suits different questions \citep{xu2024retrieval,yu2024defenserag,li2024selfroute}, and long-context benchmarks show that models use long inputs unevenly \citep{bai2024longbench,modarressi2025nolima}. Iterative and agentic search lets a model issue new queries \citep{jiang2023flare,trivedi2023ircot,jeong2024adaptiverag,asai2024selfrag,yan2024crag,yao2023react,jin2025searchr1}, but with few exceptions its extra cost is not weighed against the gain \citep{ferrazzi2026agentic,sharify2026rerank}. Benchmarks measure retrieval, question answering and RAG, mostly on English text \citep{thakur2021beir,muennighoff2023mteb,kwiatkowski2019nq,joshi2017triviaqa,petroni2021kilt,chen2024rgb,yang2024crag}, and multi-hop sets test reasoning across passages or documents \citep{yang2018hotpotqa,ho2020twowiki,trivedi2022musique,tang2024multihoprag,krishna2024frames,wang2024loong}. Answers are usually scored by language-model judges \citep{zheng2023judge,liu2023geval,kim2024prometheus,es2024ragas,saadfalcon2024ares}, whose verdicts vary between runs and differ from those of people \citep{faggioli2023perspectives,haldar2025roulette,thakur2025judging}. Small test sets make rankings fragile \citep{card2020power,sakai2014reform,sakai2016topicset}. Preregistration has been proposed to limit analysis choices made after the results are seen \citep{vanmiltenburg2021prereg}, and the amount of tuning behind a result should be reported \citep{dodge2019showyourwork}.

Consequently, a gap remains: when the evidence budget is fixed, which retrieval, selection and search choices decide whether the evidence in natural Persian and English documents reaches the reader? The study also examines how far the resulting rankings can be trusted. In order to answer the question, this paper presents a comparison of RAG setups in which the retrieval and selection stages vary while the reader and the answer prompt stay fixed. The comparison covers four whole pipelines, single-component variations, iterative and long-context alternatives, and the two languages. Only part of this design was fixed in advance: a confirmatory test of one proposed pipeline against two baselines, with decision rules that were frozen before any test answer was seen. All comparisons of components, budgets, search strategies and languages are post hoc, and they are labelled as such. The study also records, for each question, whether the evidence pack holds a complete set of supporting passages. A loss in retrieval can then be told from a loss in packing, which is the choice of what fits into the budget. Finally, it checks the rankings against repeated runs of the reader, two judges and a partial human review. One of the four pipelines, called Lex-RAG, is a selection-based design proposed in this work, and it tests whether one language-model call that selects evidence improves a hybrid pipeline.

In order to address this gap, this paper makes the following contributions:
\begin{itemize}
\item \textbf{Comparison under an evidence budget:} A comparison of whole pipelines and single components at fixed evidence budgets, with a frozen confirmatory test of the proposed pipeline and post hoc comparisons of components, budgets and search strategies that are labelled as such.
\item \textbf{A Persian and English benchmark with evidence bundles:} A benchmark of natural documents in three genres and two languages, with three tiers of evidence spread and annotated bundles of supporting passages that separate losses in retrieval from losses in packing.
\item \textbf{Diagnosis, reliability checks and guidance:} An account of where answers are lost, a check of how far the rankings survive repeated reader runs, a second judge and a different review design, and recommendations with a stated confidence.
\end{itemize}

Table~\ref{tab:related} compares the closest prior studies and benchmarks. Its columns are the aspects that this paper combines: languages, natural documents, a varied evidence budget, reranking or selection, iterative search, cross-lingual queries, human review and a frozen test. The marks follow what each cited work reports. In the row of this study, a superscript shows whether a check was fixed before the test (c) or is post hoc (p). The table shows that prior work covers these aspects separately and rarely places them in one design.
% src: marks re-checked 2026-11-18 against each paper (arXiv HTML/PDF; see the % verified lines); a dash = not stated. This study = Sects. 2-3.
\providecommand{\crossmark}{\ensuremath{\times}}
\begin{table}[!t]
\caption{Comparison of related studies and benchmarks by what each covers. A check mark means covered and a cross means not covered}\label{tab:related}
\centering\scriptsize\setlength{\tabcolsep}{2pt}
\begin{tabular*}{\textwidth}{@{\extracolsep{\fill}}L{0.15\textwidth}L{0.11\textwidth}cccccccc@{}}
\toprule
Study & Languages & \shortstack{Natural\\documents} & \shortstack{Budget\\varied} & \shortstack{Rerank or\\select} & \shortstack{Iterative\\search} & \shortstack{Cross-\\lingual} & \shortstack{Human\\review} & \shortstack{Frozen\\test}\\
\midrule
% verified: arXiv 2407.01219 (HTML). Corpus = English Wikipedia + medical text, MS MARCO/TREC DL passages (lyft 10-K pages only for the chunk-size test) -> natural x; no top-k/token sweep (chunk size and compression ratio 0.4 compared, top-k unstated) -> budget x; monoT5/monoBERT/RankLLaMA/TILDEv2 rerankers + Recomp/LongLLMLingua compression -> rerank check; GPT-4 RAGAs judge, no humans; modules picked sequentially on the reported benchmarks -> frozen x.
Wang et al.~\citep{wang2024bestpractices} & English & \crossmark & \crossmark & \checkmark & \crossmark & \crossmark & \crossmark & \crossmark\\
% verified: ACL Anthology 2024.findings-emnlp.449 (PDF read 2026-11-18). KILT Wikipedia dump split into 100-word chunks (Sec. 4.5), one Wikipedia-based collection only (Limitations) -> natural x; top-50 retrieved and top-5 given to the LLM, fixed (Sec. 5) -> budget x; MiniLM, DeBERTa-v3 and BGE-M3 cross-encoders compared with and without reranking (Secs. 4.2, 5.3, Fig. 4) -> rerank check; single-pass pipelines -> iterative x; multilingual extension with MKQA and XOR-TyDi QA, queries in ar, fi, ja, ko, ru against English or multilingual Wikipedia (Sec. 5.5, Table 2) -> cross-lingual check (small extension), languages English plus those five (arXiv 2407.01102 Sect. 5.6, Table 3 header en, ar, fi, ja, ko, ru; rechecked 2026-11-18); metrics and LLM judges (LLMeval, GPT-4), no human annotation of answers -> human x; no pre-specified test, best practices derived from the experiments -> frozen x.
Rau et al.~\citep{rau2024bergen} & English and five others & \crossmark & \crossmark & \checkmark & \crossmark & \checkmark & \crossmark & \crossmark\\
% verified: arXiv 2407.16833 (HTML). English-only LongBench/InfiniteBench tasks, avg 3.6k-150k words -> natural check; top-k swept 1-100 chunks (Table 7) -> budget check; dense retrievers only, Self-Route answerability gate is routing not passage selection -> rerank x; single-pass; automatic metrics only (manual failure inspection, not answer judging); no pre-specified test.
Li et al.~\citep{li2024selfroute} & English & \checkmark & \checkmark & \crossmark & \crossmark & \crossmark & \crossmark & \crossmark\\
% verified: arXiv 2410.05983 (HTML). 2018 Wikipedia dump, <100-word chunks (App. F.3) -> natural x; number of passages varied (Sec. 3.1) -> budget check; only positional 'retrieval reordering' by retriever score, no reranker or LLM filter -> rerank x; no iterative, cross-lingual or human evaluation; no frozen test (language not named; NQ/TriviaQA etc. are English).
Jin et al.~\citep{jin2025longcontext} & English & \crossmark & \checkmark & \crossmark & \crossmark & \crossmark & \crossmark & \crossmark\\
% verified: arXiv 2607.00725v2 (HTML). HotpotQA/2Wiki/MuSiQue paragraphs (+ RAGBench CovidQA, ExpertQA), English bge-small-en -> natural x; reader budget B in {96,128,160,224} tokens -> budget check; top-k, MMR, submodular packer, LLMLingua-2 -> select check; EM/F1 only, no humans; heuristic-edge conditions 'identified post hoc', asks for a pre-registered replication -> frozen x.
Bala~\citep{bala2026recall} & English & \crossmark & \checkmark & \checkmark & \crossmark & \crossmark & \crossmark & \crossmark\\
% verified: arXiv 2605.27789v1 (HTML). 687 English arXiv papers (CS/ML, Materials Science), chunked -> natural check; 2000-token budget fixed for all arms, never varied -> budget x; GADMEC evolutionary evidence selector vs top-k/MMR/BM25 -> select check; single-pass; LLM judges only ('does not replace human gold annotation'); primary registration locks four hypotheses, cluster bootstrap, Bonferroni 0.0125 -> frozen check (caveats: synthetic questions, held-out dev set not stated, logged deviations: judge 4.6->4.7, n 100->200).
Chac{\'o}n Sartori and Garc{\'i}a~\citep{chacon2026fixedbudget} & English & \checkmark & \crossmark & \checkmark & \crossmark & \crossmark & \crossmark & \checkmark\\
% verified: arXiv 2601.07711 (HTML). NQ, FEVER (Wikipedia), CQADupStack-English, FiQA (short forum posts) -> natural x; fixed number of retrieved chunks, no sweep -> budget x; ELECTRA reranker as a tested dimension (Sec. 3.4) -> rerank check; agentic RAG re-retrieves -> iterative check; manual annotation of 312 answer pairs (5% of CQAD test), 65.4% human-model agreement -> human check; no pre-registration.
Ferrazzi et al.~\citep{ferrazzi2026agentic} & English & \crossmark & \crossmark & \checkmark & \checkmark & \crossmark & \checkmark & \crossmark\\
% verified: arXiv 2104.08663 (PDF). 18 English datasets; mixed corpora, incl. long news articles (TREC-NEWS 635, Robust04 466 words avg) besides Wikipedia/short passages -> natural check (mixed); retrieval-only benchmark, no reader -> budget x; BM25+CE reranks top-100 -> rerank check; the authors' manual annotation (980 TREC-COVID query-doc pairs) judges relevance, not answers -> human x; zero-shot test sets but no pre-fixed decision rules -> frozen x.
BEIR~\citep{thakur2021beir} & English & \checkmark & \crossmark & \checkmark & \crossmark & \crossmark & \crossmark & \crossmark\\
% verified: arXiv 2406.04744 (HTML). Up to 50 full HTML web pages per question + mock KG APIs -> natural check; tasks give 5 vs 50 pages but the reader's context is truncated to a fixed 1K-4K-token window per model -> budget x; no reranker/filter evaluated; single-pass baselines; two human graders (94% agreement) score answers -> human check; private split held out for KDD Cup winners, but no pre-fixed decision rules for the paper's own comparisons -> frozen x.
CRAG~\citep{yang2024crag} & English & \checkmark & \crossmark & \crossmark & \crossmark & \crossmark & \checkmark & \crossmark\\
% verified: arXiv 2407.01463 (HTML) / ACL Anthology 2024.knowllm-1.15. 13 languages; KILT/Wikipedia 100-word chunks -> natural x; fixed 50 retrieved -> top-5 -> budget x; BGE-m3 reranker always on, retrievers compared but no reranker/selection comparison -> rerank x; user-language queries vs English Wikipedia -> cross-lingual check; manual error analysis of 50 MKQA answers each for ru/zh/fr (Table 8) -> human check (small); prompts chosen after preliminary runs -> frozen x.
Chirkova et al.~\citep{chirkova2024mrag} & Multiple & \crossmark & \crossmark & \crossmark & \crossmark & \checkmark & \checkmark & \crossmark\\
% verified: arXiv 2410.21970 (HTML). 8 languages; 197 Wikipedia-derived documents with fictionalised entities/dates -> natural x; top-5 200-token chunks fixed (k varies only the language mix of documents) -> budget x; no reranker or selection step (reader's language preference among conflicting docs is the measured outcome) -> rerank x; single-pass; cross-language document/query pairs -> cross-lingual check; experts reviewed the data, answers scored by 3-gram recall / GPT-4 -> human x; no held-out protocol.
Wu et al.~\citep{wu2024multilingual} & Multiple & \crossmark & \crossmark & \crossmark & \crossmark & \checkmark & \crossmark & \crossmark\\
% verified: arXiv 2507.07543 (HTML). UAE legislation (390 laws, ~1.5M words) and MoFA travel pages, split into 100-token passages -> natural check; fixed top-20 -> rerank -> top-5 -> budget x; one BGE-v2-M3 reranker, never compared with alternatives or no-rerank -> rerank x; single-pass; query and document language chosen independently -> cross-lingual check; native speakers checked the LLM judge's answer labels (200 Wikipedia QA pairs, 95%/98% agreement) -> human check; no pre-specified test.
Amiraz et al.~\citep{amiraz2025crosslingual} & Arabic, English & \checkmark & \crossmark & \crossmark & \crossmark & \checkmark & \checkmark & \crossmark\\
% verified: arXiv 2501.04858 (PDF). PQuAD (Wikipedia) + a Persian physical-education textbook + the Fundamental Transformation Document of Education (Sec. 5.1) -> natural check (2 of 3 corpora); chunk size 512/1024/2048 at similarity top-k 5 (Table 4, Sec. 5.3.3) -> budget check; LLM document-summary index that first judges summaries for relevance, with/without (Sec. 5.3.4, Table 9) -> select check; Persian only; RAGAS/MCQ accuracy, no human judging; no held-out protocol.
Hosseinbeigi et al.~\citep{hosseinbeigi2025persianrag} & Persian & \checkmark & \checkmark & \checkmark & \crossmark & \crossmark & \crossmark & \crossmark\\
\midrule
% Item 3 (REVIEW_J1 M2): only the Lex-RAG gate and the two primary contrasts were fixed in advance (c); the rest is post hoc (p).
% Budget: 6,000 and 15,500 characters frozen, 3,000, 10,000 and 24,000 post hoc. Rerank/select: Lex-RAG and HR frozen, component variants post hoc.
% Iterative search and cross-lingual (machine-translated questions): post hoc single runs. Human review: the author, 30 questions, after the machine results.
\textbf{This study (proposed)} & Persian, English & \checkmark & \checkmark\textsuperscript{c,p} & \checkmark\textsuperscript{c,p} & \checkmark\textsuperscript{p} & \checkmark\textsuperscript{p} & \checkmark\textsuperscript{p} & \checkmark\textsuperscript{c}\\
\bottomrule
\end{tabular*}
\par\smallskip
{\scriptsize\emph{Note.} For this study, c marks a check that was fixed before the test and p a post hoc check. The frozen test covers one proposed pipeline against two baselines, with a private code commit as its record. Human review is by the author on 30 questions. Cross-lingual queries are machine-translated.\par}
\end{table}


The remainder of this paper is organised as follows. Section~\ref{sec:methods} presents the methodology, with Sect.~\ref{sec:formulation} formulating the task, Sect.~\ref{sec:benchmark} describing the benchmark, Sect.~\ref{sec:systems} the pipelines and component variations, Sect.~\ref{sec:readers} the readers and judges, and Sect.~\ref{sec:protocol} the experimental setup. Section~\ref{sec:results} reports the experimental evaluation. It starts with the confirmatory test in Sect.~\ref{sec:frozen} and examines retrieval components and where answers are lost in Sect.~\ref{sec:components}. It then studies the evidence budget with iterative and long-context retrieval in Sect.~\ref{sec:evidence}, compares Persian and English in Sect.~\ref{sec:language}, and assesses the reliability of the rankings, the practical guidance and the limitations in Sect.~\ref{sec:stability}. Finally, the paper concludes in Sect.~\ref{sec:conclusion} with a summary of the findings and directions for future work.
